% Figura corregida: radar conceptual de dimensiones de evaluación.
\begin{figure}[H]
\centering
\resizebox{0.78\textwidth}{!}{%
\begin{tikzpicture}[font=\sffamily]
\tikzset{axis/.style={draw=institucionalBlue, line width=0.55pt}, poly/.style={draw=institucionalRed, fill=white, line width=1.0pt}, lab/.style={font=\scriptsize, align=center}}
\foreach \r in {1,2,3,4}{\draw[axis] (0,0) circle (\r);} 
\foreach \ang/\txt in {90/Planeación,30/Ejecución,-30/Cierre\\administrativo,-90/Trazabilidad,-150/Usabilidad,150/Seguridad}{
    \draw[axis] (0,0) -- (\ang:4.15);
    \node[lab] at (\ang:4.85) {\txt};
}
\foreach \r/\txt in {1/básico,2/parcial,3/alto,4/completo}{\node[font=\tiny] at (8:\r) {\txt};}
\draw[poly]
    (90:3.0) -- (30:3.2) -- (-30:2.5) -- (-90:3.1) -- (-150:2.7) -- (150:3.3) -- cycle;
\node[font=\small\bfseries, text=institucionalBlue] at (0,5.55) {Dimensiones de evaluación funcional del aplicativo};
\node[figNote, text width=4.6cm] at (5.9,-0.1) {Representación cualitativa. Sustituir por datos cuantitativos solo si se anexan mediciones verificables.};
\end{tikzpicture}%
}
\caption[Radar cualitativo de dimensiones de evaluación]{Radar cualitativo de dimensiones de evaluación del aplicativo: planeación, ejecución, cierre, trazabilidad, usabilidad y seguridad. Fuente: Elaboración propia.}
\label{fig:radar-mejoras}
\end{figure}
